\scanpage{126}
\margin{74}
\[
H^1(\mathcal E_\phi,O_{\mathcal E_\phi}[nE]_{\mathcal E_\phi})=0
\]
we obtain an exact sequence
\diagram{\begin{tikzcd}
H^0(V,O_V[nE]) \arrow[r,"r^*"] & H^0(\mathcal E_\phi,O_{\mathcal E_\phi}[nE]_{\mathcal E_\phi}) \arrow[r] & 0
\end{tikzcd}}\hfill(19)

Let us assume that there would exist a point $P\in\mathcal E_\phi$ such that in $P$ any cross-section $\phi\in H^0(V,O_V[nE])$ satisfies the condition: if $P\in U^{(i)}$, then $\phi^{(i)}(P)=0$. Since any cross-section $\phi'\in H^0(\mathcal E_\phi,O_{\mathcal E_\phi}[nE]_{\mathcal E_\phi})$ can be presented in the form $\phi'=r^*\phi$ or $\phi'^{(i)}=r\phi^{(i)}$ and $r$ in the stalk over $P$ is a homomorphism of restriction of functions; then $\phi'^{(i)}(P)=0$. But this contradicts the fact that, at $n\ge n(\mathcal E_\phi)$, the elements of the group $H^0(\mathcal E_\phi,O_{\mathcal E_\phi}[nE]_{\mathcal E_\phi})$ generate the stalk of the sheaf $O_{\mathcal E_\phi}[nE]_{\mathcal E_\phi}$ over any point $P\in\mathcal E_\phi$. The lemma is proved.\ \rule{0.5em}{1em}

Let us come back to the proof of theorem 2. Put $E'=n_{\mathcal E}E$, where $n_{\mathcal E}$ is the number to which lemma 3 refers. Let $\phi_0,\ldots,\phi_T$ be the basis (not necessarily of linearly independent vectors) of the group $H^0(V,O_V[nE'])$.

Consider a mapping $f_n:\ V\to P^T$ ($P^T$ is a T-dimensional projective space), operating according to the rule: if $P\in V$, $P\in U^{(i)}$, then
\mspage{125}
